\documentclass[a4paper]{article}

% Español (México) con XeLaTeX: fontspec para fuentes Unicode, polyglossia
% para la división silábica y los nombres del documento en la variante mexicana.
\usepackage{fontspec}
\usepackage{polyglossia}
\setdefaultlanguage[variant=mexican]{spanish}
% Los nombres chinos van como «pinyin (original)»: xeCJK da a los caracteres
% Han su propia fuente; sin ella XeLaTeX los omite con un aviso.
% FandolSong no cubre algunos caracteres de nombres (U+73FA); WenQuanYi Zen Hei sí.
\usepackage[AutoFallBack=true]{xeCJK}
\setCJKmainfont{FandolSong-Regular.otf}
\setCJKfallbackfamilyfont{\CJKrmdefault}{WenQuanYi Zen Hei}
% polyglossia no traduce los nombres de listings.
\AtBeginDocument{\providecommand{\lstlistingname}{}\renewcommand{\lstlistingname}{Listado}%
  \providecommand{\lstlistlistingname}{}\renewcommand{\lstlistlistingname}{Índice de listados}}
\usepackage{amsmath,amssymb}
\usepackage{graphicx}
\usepackage[margin=2.5cm]{geometry}
\usepackage[most]{tcolorbox}
\usepackage{etoolbox}
\usepackage{subcaption}
\usepackage{listings}
\usepackage{hyperref}
\usepackage{booktabs}
\usepackage{float}

\newtcolorbox{knowledgebox}[1]{
    enhanced, colback=blue!5!white, colframe=blue!75!black, colbacktitle=blue!75!black,
    coltitle=white, fonttitle=\bfseries, title={#1}, boxrule=1pt, sharp corners
}

\newtcolorbox{importantbox}[1]{
    enhanced, colback=yellow!10!white, colframe=yellow!80!black, colbacktitle=yellow!80!black,
    coltitle=black, fonttitle=\bfseries, title={#1}, sharp corners
}

\newtcolorbox{warningbox}[1]{
    enhanced, colback=red!5!white, colframe=red!75!black, colbacktitle=red!75!black,
    coltitle=white, fonttitle=\bfseries, title={#1}, sharp corners
}

\lstset{
    basicstyle=\ttfamily\small,
    keywordstyle=\color{blue},
    stringstyle=\color{red!60!black},
    commentstyle=\color{green!60!black},
    breaklines=true,
    frame=single,
    numbers=left,
    numberstyle=\tiny\color{gray},
    captionpos=b,
    extendedchars=true
}

\newcommand{\notetitle}{CS224R Lecture 14: Exploración y metaexploración}
\newcommand{\noteauthors}{Elaboradas a partir de materiales públicos del curso}
\newcommand{\notedate}{Primavera de 2025}
\newcommand{\videochannel}{Stanford Online}
\newcommand{\videopublishdate}{2025}
\newcommand{\videoduration}{Aproximadamente 80 minutos}
\newcommand{\videourl}{https://cs224r.stanford.edu/}
\newcommand{\videocoverpath}{cover.jpg}
\newcommand{\repourl}{https://github.com/hqhq1025/ai-course-notes}


% Footer with repo info
\usepackage{fancyhdr}
\pagestyle{fancy}
\fancyhf{}
\fancyfoot[C]{\thepage}
\fancyfoot[R]{\footnotesize\color{gray}\href{\repourl}{AI Course Notes}}
\renewcommand{\headrulewidth}{0pt}
\renewcommand{\footrulewidth}{0pt}

\begin{document}
\begin{titlepage}
\centering
{\Large Notas del curso\par}
\vspace{1.2cm}
{\huge\bfseries \notetitle\par}
\vspace{0.8cm}
{\large \noteauthors\par}
\vspace{0.3cm}
{\large \notedate\par}
\vspace{1.1cm}

\ifdefempty{\videocoverpath}{
{\small No se proporcionó imagen de portada.\par}
}{
\includegraphics[width=0.82\textwidth,height=0.45\textheight,keepaspectratio]{\videocoverpath}\par
}

\vfill
\begin{tcolorbox}[width=0.9\textwidth, colback=black!2!white, colframe=black!60, sharp corners]
\textbf{Autor o canal del video}: \videochannel\par
\textbf{Fecha de publicación}: \videopublishdate\par
\textbf{Duración del video}: \videoduration\par
\textbf{Ponente}: Chelsea Finn\par
\textbf{Página del curso}: \href{https://cs224r.stanford.edu/}{cs224r.stanford.edu}\par
\textbf{Página del proyecto}: \href{\repourl}{\nolinkurl{\repourl}}\par
\end{tcolorbox}
\end{titlepage}

\tableofcontents
\newpage

%% --- Inicio del contenido --- %%
\section{La naturaleza del problema de la exploración}

\subsection{Por qué importa la exploración}

En el aprendizaje por refuerzo, el agente debe encontrar un equilibrio entre la \textbf{explotación} (exploitation, elegir la mejor acción conocida) y la \textbf{exploración} (exploration, probar acciones que podrían ser mejores).

\begin{figure}[H]
\centering
\includegraphics[width=0.9\textwidth]{slides-images/slide-04.jpg}
\caption{Ilustración intuitiva de la dificultad de la exploración: algunas tareas son sencillas para una persona, pero casi imposibles para un RL que parte de cero}
\end{figure}

\begin{importantbox}{El dilema central de la exploración}
Si siempre elegimos la acción con el valor estimado más alto en ese momento (política voraz), quizá nunca descubramos la política verdaderamente óptima. Por ejemplo, en el bandido multibrazo, si al accionar un brazo por primera vez se obtiene por casualidad una recompensa positiva, la política voraz seguirá accionando ese mismo brazo, aunque otros brazos tengan una recompensa esperada mayor.
\end{importantbox}

\subsection{Por qué Montezuma's Revenge se volvió una referencia clásica de exploración}

La clase empieza ilustrando el problema con un juego de recompensas escasas como Montezuma's Revenge: obtener la llave da recompensa y abrir la puerta da recompensa, pero morir a manos de la skull no necesariamente produce de inmediato una retroalimentación negativa lo bastante clara, y los comportamientos decisivos para completar la tarea están muy separados entre sí. Una persona sabe qué hacer no porque vea la reward, sino porque entiende la semántica de los sprites y la estructura de la tarea.

\begin{figure}[H]
\centering
\includegraphics[width=0.9\textwidth]{slides-images/slide-05.jpg}
\caption{Montezuma's Revenge: las recompensas escasas, las cadenas largas de comportamiento y la falta de comprensión semántica dificultan la exploración}
\end{figure}

\subsection{Ponerse en el lugar del algoritmo}

El ponente lleva la analogía más lejos con el juego de cartas Mao: solo sabes que «romper las reglas se castiga», pero las reglas mismas solo pueden descubrirse por ensayo y error, y no siempre son intuitivas. Esto describe con precisión la dificultad esencial de las tareas de horizonte largo con recompensas escasas.

\begin{figure}[H]
\centering
\includegraphics[width=0.9\textwidth]{slides-images/slide-06.jpg}
\caption{Analogía de Mao: las reglas solo se descubren por ensayo y error; cuanto más larga es la tarea y más ocultas están las reglas, más difícil resulta explorar}
\end{figure}

\subsection{Exploración y explotación son en realidad el mismo problema}

La clase subraya que dos formulaciones habituales son, en esencia, equivalentes:
\begin{itemize}
    \item «¿Cómo descubrir políticas de alta recompensa a las que solo se llega mediante comportamientos complejos y de cadena larga?»
    \item «¿Cuándo conviene seguir haciendo lo mejor que ya se conoce y cuándo conviene probar un comportamiento nuevo?»
\end{itemize}
La primera pone el acento en el temporally extended behavior; la segunda, en el exploration-exploitation trade-off; pero ambas preguntan lo mismo: \textbf{¿estamos dispuestos a aceptar hoy una pequeña pérdida a cambio de un rendimiento posiblemente mayor mañana?}

\begin{figure}[H]
\centering
\includegraphics[width=0.9\textwidth]{slides-images/slide-07.jpg}
\caption{Restaurantes, colocación de anuncios y prospección de pozos petroleros: ejemplos cotidianos de exploración y explotación}
\end{figure}

\begin{warningbox}{La exploración no es difícil solo porque el espacio de acciones sea grande}
Lo verdaderamente complicado es que muchos comportamientos de exploración decisivos no muestran beneficio a nivel local e incluso reducen el rendimiento a corto plazo. Por eso, si la señal de entrenamiento solo atiende al desempeño inmediato, el agent tiende a verse empujado hacia políticas locales conservadoras pero subóptimas.
\end{warningbox}

\begin{figure}[H]
\centering
\includegraphics[width=0.9\textwidth]{slides-images/slide-08.jpg}
\caption{Del bandit a los MDP grandes y sin estructura: la dificultad teórica del problema de la exploración crece rápidamente}
\end{figure}

\subsection{La exploración en las aplicaciones de RL}

\begin{itemize}
    \item \textbf{Robótica}: explorar distintas estrategias de agarre y distintos patrones de movimiento
    \item \textbf{Modelos de lenguaje grandes}: en RLHF/RL for reasoning, explorar distintas rutas de razonamiento
    \item \textbf{Sistemas de recomendación}: explorar contenido que podría gustarle al usuario pero que todavía no se le ha mostrado
\end{itemize}

\subsection{Resumen de la sección}

La exploración es uno de los retos fundamentales de todo problema de RL. Sin una exploración eficaz, el agente se queda atrapado en políticas subóptimas.
\section{Métodos clásicos de exploración}

\subsection{El planteamiento Multi-Armed Bandit}

\begin{knowledgebox}{Multi-Armed Bandit}
El problema bandit es la versión más simplificada del problema de exploración: hay $K$ brazos (acciones) y cada brazo tiene una distribución de recompensa desconocida $p(r | a)$. No hay transiciones de estado; el objetivo es maximizar la recompensa acumulada a lo largo de $T$ rondas de interacción. El arrepentimiento (regret) se define como la diferencia respecto al brazo óptimo:
$$\text{Regret}(T) = T \cdot \mu^* - \sum_{t=1}^T r_t$$
\end{knowledgebox}

\subsection{Cómo se define una «buena exploración» en un bandit}

El escenario bandit es importante porque en él todavía puede definirse con bastante claridad la «exploración óptima». El curso adopta el regret como métrica central: compara la recompensa acumulada que realmente obtienes con el rendimiento esperado de tirar siempre del brazo óptimo.

\begin{figure}[H]
\centering
\includegraphics[width=0.9\textwidth]{slides-images/slide-09.jpg}
\caption{Medida de optimalidad en un bandit: el regret, no la recompensa de una sola ronda}
\end{figure}

\subsection{$\epsilon$-Greedy}

La estrategia de exploración más sencilla: con probabilidad $1 - \epsilon$ se elige la acción que actualmente es la mejor y con probabilidad $\epsilon$ se elige una al azar.

\begin{warningbox}{Limitaciones de $\epsilon$-Greedy}
La exploración de $\epsilon$-Greedy es \textbf{no dirigida}: explora de manera aleatoria y uniforme sobre todas las acciones, sin considerar cuáles aportan más información. Cuando el espacio de acciones es grande o el entorno es complejo, esta exploración a ciegas resulta sumamente ineficiente.
\end{warningbox}

\subsection{Upper Confidence Bound (UCB)}

\begin{importantbox}{Algoritmo UCB}
La idea central de UCB es el \textbf{optimismo} (optimism in the face of uncertainty): asigna una estimación más alta a las acciones con mayor incertidumbre y, de ese modo, fomenta la exploración.

$$a_t = \arg\max_a \left[\hat{\mu}_a + c \sqrt{\frac{\ln t}{N_a}}\right]$$

donde $\hat{\mu}_a$ es la media empírica de la acción $a$ y $N_a$ es el número de veces que se ha elegido la acción $a$. El segundo término es una bonificación por incertidumbre: cuantas menos veces se ha elegido una acción, mayor es su incertidumbre y con mayor prioridad se explora.
\end{importantbox}

\subsection{Thompson Sampling}

\begin{importantbox}{Thompson Sampling}
Thompson Sampling adopta un enfoque bayesiano:
\begin{enumerate}
    \item Para la distribución de recompensa de cada acción se mantiene una distribución posterior $p(\mu_a | \text{data})$
    \item En cada ronda se muestrea de la posterior: $\tilde{\mu}_a \sim p(\mu_a | \text{data})$
    \item Se elige la acción con el valor muestreado más alto: $a_t = \arg\max_a \tilde{\mu}_a$
    \item Tras obtener la recompensa, se actualiza la posterior
\end{enumerate}
En la práctica, Thompson Sampling suele tener un desempeño excelente y, además, cuenta con un regret bound teóricamente óptimo.
\end{importantbox}

\begin{figure}[H]
\centering
\includegraphics[width=0.9\textwidth]{slides-images/slide-10.jpg}
\caption{Las dos líneas principales de la exploración en bandits: optimism under uncertainty y probability matching}
\end{figure}

\subsection{Por qué la teoría de bandits sigue siendo el punto de partida de la investigación sobre exploración}

El bandit parece demasiado simple, pero sigue siendo la vía de entrada clave para entender el problema de la exploración, porque es uno de los pocos planteamientos en los que todavía podemos ofrecer garantías de regret y discutir la «exploración óptima». Al pasar a MDP de horizonte largo, parcialmente observables y con estados de alta dimensión, muchas de estas conclusiones dejan de cumplirse directamente; por eso UCB y Thompson Sampling funcionan más como principios de diseño que se conservan que como plantillas de algoritmos que puedan copiarse tal cual.

\begin{knowledgebox}{Tres principios de diseño que se aprenden del bandit}
\begin{enumerate}
    \item La incertidumbre debe formar parte de la decisión; no basta con mirar la media empírica.
    \item La calidad de la exploración se juzga por el costo acumulado a largo plazo, no por el desempeño en un solo paso.
    \item En cuanto la tarea pasa de un 1-step bandit a una toma de decisiones de horizonte largo, la dificultad de la exploración aumenta drásticamente.
\end{enumerate}
\end{knowledgebox}

\subsection{Resumen de la sección}

Los métodos clásicos de exploración van desde el sencillo $\epsilon$-greedy hasta los métodos con fundamento de UCB y Thompson Sampling; la diferencia central entre ellos está en cómo aprovechan la incertidumbre para guiar la exploración.
\section{Exploración en RL profundo}

\subsection{De Bandit a MDP}

En un MDP, el problema de exploración es más complejo:
\begin{itemize}
    \item Las acciones no solo afectan la recompensa inmediata, sino también los estados futuros
    \item Es necesario explorar \textbf{pares estado--acción} y no solo acciones
    \item Algunos estados son difíciles de alcanzar en sí mismos y requieren una exploración dirigida durante mucho tiempo
\end{itemize}

\subsection{Exploración basada en conteo}

\begin{importantbox}{Count-Based Exploration}
Idea básica: otorgar una recompensa adicional de exploración a los pares estado--acción que se han visitado pocas veces:
$$r_{\text{explore}}(s, a) = \frac{\beta}{\sqrt{N(s, a)}}$$
En RL profundo, como el espacio de estados es continuo y de alta dimensión, no es posible contar directamente. Los métodos de aproximación más comunes incluyen:
\begin{itemize}
    \item Conteo con hash: mapear los estados a un espacio discreto mediante una función hash
    \item Modelo de densidad: entrenar un modelo de densidad $\hat{p}(s)$ y usar $-\log \hat{p}(s)$ como medida de novedad
    \item RND (Random Network Distillation): usar el error de predicción respecto de una red aleatoria como medida de novedad
\end{itemize}
\end{importantbox}

\subsection{Exploración basada en curiosidad}

Otro enfoque: recompensa de exploración = error de predicción del modelo. Si el error de predicción para un estado es grande, eso indica que el estado es «novedoso» y vale la pena seguir explorándolo.

\begin{warningbox}{La trampa de la curiosidad}
Los métodos de exploración basados en el error de predicción tienen un modo de falla conocido: el \textbf{problema de la televisión con ruido} (noisy TV problem). Si el entorno contiene aleatoriedad intrínsecamente impredecible (como la estática en la pantalla de un televisor), el error de predicción del modelo nunca disminuye, lo que hace que el agente quede «atraído» por esa aleatoriedad sin sentido y no pueda continuar con una exploración significativa.
\end{warningbox}

\subsection{Por qué los robots y los LLM rara vez exploran desde cero}

La conclusión que el curso da en esta parte es bastante pragmática: en MDP enormes como el control de robots y los LLM, realizar una exploración eficaz desde cero suele ser computacionalmente inasumible. En la industria y en la investigación de frontera, las estrategias más comunes son:
\begin{itemize}
    \item usar demonstrations o un base model preentrenado para reducir el espacio de búsqueda;
    \item proporcionar shaped rewards siempre que sea posible;
    \item reservar la exploración realmente difícil para subproblemas más pequeños y mejor estructurados.
\end{itemize}

\begin{figure}[H]
\centering
\includegraphics[width=0.9\textwidth]{slides-images/slide-13.jpg}
\caption{En robótica y en los LLM, normalmente se recurre al preentrenamiento, a las demostraciones y a shaped reward, en lugar de explorar desde cero}
\end{figure}

\subsection{Resumen de la sección}

La exploración en RL profundo requiere manejar medidas de novedad en espacios continuos de alta dimensión; los métodos count-based y los métodos basados en curiosidad tienen cada uno sus ventajas y limitaciones.
\section{Metaexploración: aprender a explorar}

\subsection{Qué es la metaexploración}

\begin{importantbox}{Meta-Exploration}
La metaexploración (meta-exploration) aplica las ideas del metaaprendizaje al problema de la exploración: en lugar de diseñar a mano una política de exploración, se \textbf{aprende una política de exploración}.

En el marco de meta-RL:
\begin{itemize}
    \item El nivel externo (meta) aprende a explorar de manera eficaz en tareas nuevas
    \item El nivel interno (task) usa los datos reunidos durante la exploración para adaptarse rápidamente a la tarea nueva
\end{itemize}
\end{importantbox}

\subsection{El reto de la exploración en meta-RL}

En meta-RL, las primeras interacciones del agente con cada tarea nueva son decisivas: determinan si el agente puede inferir rápidamente las características de la tarea. Por ello, la metaexploración se ocupa de \textbf{cómo reunir los datos más informativos con un número limitado de interacciones}.

\begin{knowledgebox}{Diferencia con la exploración clásica}
La exploración clásica se ocupa del equilibrio exploration-exploitation dentro de una sola tarea. La metaexploración se ocupa de aprender una política de exploración que se aplique entre tareas: entrenar una política capaz de explorar con eficacia en tareas nuevas. Se trata de un problema de optimización de nivel superior.
\end{knowledgebox}

\begin{figure}[H]
\centering
\includegraphics[width=0.9\textwidth]{slides-images/slide-15.jpg}
\caption{Cuando aún hay que explorar en el momento de prueba: buscar ingredientes en una cocina o elegir una estrategia de razonamiento en un LLM son exploraciones al estilo de meta-RL}
\end{figure}

\subsection{El problema de acoplamiento: por qué suele fallar aprender la exploración de extremo a extremo}

Si tanto la exploración como la ejecución se entrenan únicamente con la recompensa final de la tarea, surge el típico dilema chicken-and-egg: si la exploración es deficiente, la fase de ejecución no obtiene una buena reward; y si la fase de ejecución fracasa siempre, la fase de exploración tampoco recibe una señal de entrenamiento lo bastante clara.

\begin{figure}[H]
\centering
\includegraphics[width=0.9\textwidth]{slides-images/slide-16.jpg}
\caption{El coupling problem de la meta-learning with exploration de extremo a extremo}
\end{figure}

\subsection{DREAM y los métodos de muestreo posterior}

El ponente explicó cómo lograr una exploración eficaz en meta-RL:
\begin{itemize}
    \item Mantener una creencia (belief) sobre la identidad de la tarea
    \item Usar las ideas de Thompson Sampling para guiar la exploración
    \item Entrenar con RL en el nivel externo para que la política de exploración interna maximice la ganancia de información
\end{itemize}

\subsection{Estrategias alternativas I: muestreo posterior, recompensas intrínsecas y predicción de la tarea}

La clase repasó primero varias líneas de trabajo existentes: PEARL usa posterior sampling, MAME usa intrinsic rewards y MetaCURE usa task dynamics/reward prediction. Tienen en común que \textbf{ya no dejan la exploración por completo en manos de la señal de reward end-to-end, sino que introducen explícitamente una estructura de inferencia de la tarea.}

\begin{figure}[H]
\centering
\includegraphics[width=0.9\textwidth]{slides-images/slide-17.jpg}
\caption{Alternative strategies: posterior sampling, intrinsic rewards, task dynamics prediction}
\end{figure}

\subsection{DREAM: desacoplar explícitamente la exploración y la ejecución}

El verdadero protagonista de esta clase es DREAM. Su idea central no es aprender directamente una política de caja negra que se encargue a la vez de explorar y de ejecutar, sino proceder en dos pasos:
\begin{enumerate}
    \item Aprender a ejecutar y, al mismo tiempo, identificar una representación de cuello de botella de información que capture lo que realmente distingue a una tarea de otra.
    \item Aprender después una política de exploración dedicada específicamente a recuperar esa información relevante para la tarea.
\end{enumerate}

\begin{figure}[H]
\centering
\includegraphics[width=0.9\textwidth]{slides-images/slide-20.jpg}
\caption{La idea de desacoplamiento de DREAM: primero identificar la información relevante para la tarea y luego entrenar la exploración para recuperarla}
\end{figure}

\begin{knowledgebox}{Por qué el cuello de botella de información es necesario}
Si la representación de la tarea no se comprime, la política de exploración se desvía con facilidad hacia información decorativa ajena a la tarea. DREAM recurre a una bottlenecked representation que conserva solo las variables que de verdad influyen en la política de ejecución, como wall color o la ubicación de los ingredientes, de modo que la recompensa de la política de exploración se concentre en «encontrar la información correcta».
\end{knowledgebox}

\begin{figure}[H]
\centering
\includegraphics[width=0.9\textwidth]{slides-images/slide-22.jpg}
\caption{Cómo se entrena DREAM: la política de ejecución aprende a depender de la bottlenecked task representation y la política de exploración aprende a recuperarla}
\end{figure}

\subsection{Convertir la recuperación de información directamente en recompensa de exploración}

La clase formula además la recompensa de exploración con toda claridad: los datos $D_{\text{train}}$ que reúne la política de exploración deben ayudar lo más posible a predecir la representación de la tarea $z_i$. Así, la reward ya no depende de si la tarea final tuvo éxito, sino de «cuánto más sabemos sobre la tarea actual gracias a esta ronda de exploración».

\begin{figure}[H]
\centering
\includegraphics[width=0.9\textwidth]{slides-images/slide-23.jpg}
\caption{La recompensa central de DREAM: per-step information gain, en lugar de la reward de la tarea final}
\end{figure}

\subsection{Por qué conviene modelar por separado el cuello de botella de información}

Si en la representación de la tarea se meten demasiados detalles ajenos a la ejecución, la política de exploración intentará recuperar información equivocada, como decoraciones, texturas o fondos irrelevantes. Por eso la clase repasa expresamente el variational information bottleneck: mediante ruido y restricciones de regularización, la representación se comprime hasta conservar solo la parte de las diferencias entre tareas que la ejecución realmente necesita.

\begin{figure}[H]
\centering
\includegraphics[width=0.9\textwidth]{slides-images/slide-21.jpg}
\caption{Variational information bottleneck: comprimir la representación de la tarea y conservar solo la información relevante para la ejecución}
\end{figure}

\subsection{Resultados teóricos y experimentales}

Las conclusiones que presentan las slides son muy representativas: en un bandit-like setting, la complejidad de muestra de DREAM es mejor que la de RL$^2$; en una tarea de recompensa dispersa como 3D visual navigation, los métodos de extremo a extremo quedan claramente rezagados por el problema de coupling, mientras que DREAM alcanza un retorno mayor gracias a una exploración casi óptima.

\begin{figure}[H]
\centering
\includegraphics[width=0.9\textwidth]{slides-images/slide-24.jpg}
\caption{Análisis teórico: en ciertos escenarios, DREAM conserva las propiedades de exploración óptima y a la vez mejora notablemente la complejidad de muestra}
\end{figure}

\begin{figure}[H]
\centering
\includegraphics[width=0.9\textwidth]{slides-images/slide-25.jpg}
\caption{Tarea de navegación visual 3D con recompensa dispersa: primero leer el letrero y después rodear los obstáculos hasta llegar al objetivo correcto}
\end{figure}

\begin{figure}[H]
\centering
\includegraphics[width=0.9\textwidth]{slides-images/slide-26.jpg}
\caption{Resultados experimentales en navegación visual 3D con recompensa dispersa: DREAM se acerca al óptimo y los métodos de extremo a extremo se ven claramente afectados por el coupling}
\end{figure}

\begin{figure}[H]
\centering
\includegraphics[width=0.9\textwidth]{slides-images/slide-27.jpg}
\caption{Síntesis de los tres tipos de enfoque: de extremo a extremo, estrategias alternativas y desacoplamiento de exploración y ejecución}
\end{figure}

\subsection{Resumen de la sección}

La metaexploración convierte la propia política de exploración en objeto de aprendizaje y, mediante el entrenamiento en múltiples tareas, aprende conductas de exploración eficientes.
\section{Síntesis y ampliación}

\begin{enumerate}
    \item La exploración es el desafío central del RL; de $\epsilon$-greedy a UCB/Thompson Sampling se forma un conjunto de herramientas cada vez más refinado
    \item En el RL profundo es necesario aproximar la novedad mediante modelos de densidad o errores de predicción de redes
    \item Los métodos basados en la curiosidad deben tener cuidado con el problema de la televisión ruidosa
    \item La metaexploración trata el propio «cómo explorar» como una capacidad que puede aprenderse
    \item En aplicaciones reales como la robótica y los LLM, una exploración eficaz determina directamente la eficiencia del aprendizaje
\end{enumerate}

\subsection{Aplicación ampliada: llevar la meta-exploration a la enseñanza de la computación}

Al final, la clase también mostró una dirección poco común pero muy sugerente: usar las estrategias de exploración del meta-RL para la calificación automática y la retroalimentación sobre los programas de los estudiantes. El sistema aprende primero «qué pruebas conviene ejecutar y qué señales de comportamiento aportan más información», y después encuentra los bugs y genera sugerencias de calificación con mayor rapidez.

\newpage
\begin{figure}[H]
\centering
\includegraphics[width=0.9\textwidth]{slides-images/slide-30.jpg}
\caption{Aplicación de la meta-exploration a la retroalimentación sobre programas: explorar qué interacciones exponen mejor los problemas del código de los estudiantes}
\end{figure}

\subsection{En qué se distingue la «exploración» en el contexto educativo del RL tradicional}

Aquí la exploración ya no consiste en controlar un robot o jugar un videojuego, sino en decidir: \textbf{para determinar dónde falla el programa de un estudiante, qué interacciones debe ejecutar primero el sistema, qué señales debe observar y qué fragmentos de video debe mostrar.} El valor de este tipo de exploración está en ahorrar tiempo a los asistentes de docencia y, al mismo tiempo, en hacer que la retroalimentación sea más focalizada y más accionable.

\begin{figure}[H]
\centering
\includegraphics[width=0.9\textwidth]{slides-images/slide-31.jpg}
\caption{Learned exploration en la tarea Bounce: interacciones distintas exponen categorías de errores distintas}
\end{figure}

\begin{table}[H]
\centering
\begin{tabular}{p{0.2\textwidth} p{0.28\textwidth} p{0.28\textwidth}}
\toprule
Tarea educativa & Información que hay que explorar & Beneficio para el sistema \\
\midrule
Retroalimentación sobre tareas de programación interactivas & Qué escenarios de prueba provocan bugs con mayor probabilidad & Localizar más rápido el origen de los errores \\
Apoyo a la calificación de los asistentes de docencia & Qué evidencias bastan mejor para respaldar la calificación según la rubric & Reducir la revisión manual y los juicios repetidos \\
Análisis de patrones de error & Qué comportamientos de los estudiantes pertenecen al mismo tipo de fallo & Formar plantillas de retroalimentación más reutilizables \\
\bottomrule
\end{tabular}
\caption{Objetivos de recolección de información de la meta-exploration en la enseñanza de CS}
\end{table}

\begin{figure}[H]
\centering
\includegraphics[width=0.9\textwidth]{slides-images/slide-32.jpg}
\caption{Práctica de AI-assisted grading en el curso: un flujo de calificación de los asistentes de docencia más rápido y más preciso}
\end{figure}

\subsection{Tabla general de esta clase}

\begin{table}[H]
\centering
\begin{tabular}{p{0.18\textwidth} p{0.22\textwidth} p{0.22\textwidth} p{0.22\textwidth}}
\toprule
Enfoque & Idea central & Ventajas & Problemas principales \\
\midrule
$\epsilon$-greedy / UCB / Thompson & Usar la incertidumbre para guiar la exploración en los bandits & Analizable, con regret demostrable & Difícil de extender directamente a MDP de gran escala \\
Deep RL intrinsic exploration & Fomentar la exploración con la novedad o el error de predicción & Se adapta a espacios de estados continuos y de alta dimensión & Fácil de engañar con ruido y con novedad espuria \\
End-to-end meta-exploration & Aprender la exploración y la ejecución directamente a partir del reward de la tarea & Principio unificado, óptimo en teoría & El problema de coupling dificulta el entrenamiento \\
DREAM-style decoupling & Aprender primero la información clave de la tarea y después la exploración que la recupera & Más fácil de optimizar, buenos resultados en la práctica & Depende del diseño de la representación de la tarea y del cuello de botella de información \\
\bottomrule
\end{tabular}
\caption{Mapa de los métodos de exploración de la Lecture 14}
\end{table}

\begin{knowledgebox}{Lecturas adicionales}
\begin{itemize}
    \item Auer et al., ``Finite-time Analysis of the Multiarmed Bandit Problem,'' Machine Learning 2002
    \item Bellemare et al., ``Unifying Count-Based Exploration and Intrinsic Motivation,'' NeurIPS 2016
    \item Burda et al., ``Exploration by Random Network Distillation,'' ICLR 2019
    \item Pathak et al., ``Curiosity-driven Exploration by Self-Supervised Prediction,'' ICML 2017
    \item Liu et al., ``Explore then Execute: Adapting without Rewards via Factorized Meta-Reinforcement Learning,'' 2022
\end{itemize}
\end{knowledgebox}

\end{document}
